\chapter{Architecture}
\label{ch:architecture}

\jns{} is organised as three layers around one idea: run the expensive
exploration with machine-learning potentials, but only trust them where
they can be shown to be trustworthy.

\begin{center}
\begin{verbatim}
        +----------------------------------------------------------+
        |  credibility layer                                       |
        |  committee disagreement (UQ)  .  fallback to ORCA L2     |
        +----------------------------------------------------------+
                              ^
        +----------------------------------------------------------+
        |  engine layer                                            |
        |  reaction-network exploration (SCINE Chemoton)           |
        |  job execution daemon (SCINE Puffin)                     |
        |  NNP registered as a method family   <- this package     |
        +----------------------------------------------------------+
                              ^
        +----------------------------------------------------------+
        |  backend layer                                           |
        |  ML potentials (MACE, AIMNet2)  .  xTB  .  ORCA          |
        +----------------------------------------------------------+
\end{verbatim}
\end{center}

\section{The backend layer}

The backend layer provides every method in the method chain as an ASE
calculator (Chapter \ref{ch:calculators} documents each one).  All of them
follow the ASE \code{Calculator} protocol \cite{larsen2017ase}, so any of
them can be used stand-alone -- ideal for testing, benchmarking or custom
workflows -- without the engine:

\begin{itemize}
  \item \code{XTBCalculator} (L1): GFN2-xTB \cite{bannwarth2019gfn2} via the
    \code{xtb} command line, for fast pre-screening;
  \item \code{ORCACalculator} (L2/L3): r2SCAN-D4/def2-SVP
    \cite{furness2020r2scan,caldeweyher2019d4} for geometry optimisation and
    frequencies, and wB97M-V/def2-TZVP \cite{mardirossian2016wb97mv} for
    high-level single points;
  \item \code{NNPCommittee}: the three ML potentials plus the disagreement
    measure;
  \item \code{SafeNNPCalculator} (L0): the committee with the
    threshold-triggered fallback.
\end{itemize}

\section{The engine layer}

SCINE's exploration engine Chemoton \cite{unsleber2022chemoton} and its job
daemon Puffin \cite{bensberg2025puffin} never receive calculator objects.
They resolve a calculator from a \emph{method-family string} at a single
choke point, which makes external potentials injectable: \jns{} intercepts
that resolution step and returns Python calculators implementing SCINE's
\code{Calculator} interface.  Two building blocks do the work:

\begin{itemize}
  \item the \textbf{bridge} (\code{jnuscout.scine.bridge}) implements the
    full sixteen-method SCINE Python-calculator contract for any ASE
    calculator;
  \item the \textbf{injection patches} (\code{jnuscout.scine.puffin_patch},
    \code{readuct_patch}) replace the resolution functions on
    \code{scine_utilities.core} for the declared method families only,
    leaving every other method family on the stock path.
\end{itemize}

Chapter \ref{ch:scine-integration} documents the contract, the deployment
steps and the failure modes in detail.

\section{The credibility layer}

The committee's disagreement is the uncertainty score.  A threshold,
calibrated once on a reference set, decides per configuration whether the
MLP answer is used or the evaluation falls back to ORCA L2.  The mechanism
and the calibration protocol are the subject of Chapter
\ref{ch:uncertainty}.

\section{Design decisions kept on the record}

The following choices are easy to get wrong and were fixed deliberately;
they are the reason several rules in this design look asymmetric:

\begin{itemize}
  \item \textbf{Energy disagreement is computed within one model family.}
    Absolute energies across families are not comparable -- MACE reports
    atomization energies, AIMNet2 reports absolute total energies -- so
    subtracting across families is physically meaningless.  The default
    compares the two MACE models; the AIMNet2 member participates in the
    force disagreement and in the ensemble prediction.
  \item \textbf{Force disagreement is a per-atom maximum, not a mean.}  In
    reacting geometries the failure mode is typically one atom with a wrong
    force direction; averaging over atoms would dilute exactly the signal
    the fallback is supposed to catch.
  \item \textbf{Clones must be fully isolated.}  Puffin derives
    spin-shifted systems by cloning a calculator and writing new settings
    into the clone.  A shared settings object would silently modify the
    original; the bridge therefore copies settings and structure on
    \code{_clone_impl}.
  \item \textbf{The patched calculator declares a superset of properties and
    settings.}  The engine fails the whole job when a required property was
    not declared, so over-declaring is cheap and under-declaring is fatal.
\end{itemize}

\section{The lifecycle of one evaluation}

For a single geometry inside a running exploration, the flow through the
layers is:

\begin{enumerate}
  \item Puffin resolves the method family at its settings manager; the
    injection patch returns the bridge (Chapter \ref{ch:scine-integration}).
  \item The bridge converts the structure from atomic units to ASE units and
    calls the safe calculator.
  \item The safe calculator asks the committee: each member predicts energy
    and forces (in eV and eV/\AA{}); the disagreement scores \code{uq_force}
    and \code{uq_energy} are computed.
  \item If both scores are below their thresholds, the committee mean answer
    is returned.  Otherwise the evaluation falls back to ORCA L2 -- the
    fallback process is constructed lazily, only when it is first needed.
  \item The bridge converts the result back to atomic units and fills the
    \code{Results} object the engine expects.  Fallback counts
    (\code{n_calls}, \code{n_fallback}) accumulate for auditing.
\end{enumerate}

\begin{notebox}[Where the numbers live]
Every report of ``how often did we fall back'' in this manual follows this
lifecycle: it is a property of the calculator instance over the
configurations it was asked to evaluate, not a property of the potential
alone.
\end{notebox}
